\documentclass[11pt]{article}

\usepackage[margin=1in]{geometry}
\usepackage{amsmath, amssymb, amsthm}
\usepackage{graphicx}
\usepackage{booktabs}
\usepackage{xcolor}
\usepackage{hyperref}
\usepackage{caption}
\usepackage{subcaption}
\usepackage{bm}

\hypersetup{colorlinks=true, linkcolor=blue, citecolor=blue, urlcolor=blue}

% ── Notation shorthands ───────────────────────────────────────────────────────
\newcommand{\xvec}{\mathbf{x}}
\newcommand{\uvec}{\mathbf{u}}
\newcommand{\pvec}{\mathbf{p}}
\newcommand{\yvec}{\mathbf{y}}
\newcommand{\Xivl}{\mathbf{X}}
\newcommand{\Pivl}{\mathbf{P}}
\newcommand{\R}{\mathbb{R}}
\newcommand{\II}{\mathcal{I}}
\newcommand{\Lsep}{\mathcal{L}_{\rm sep}}
\newcommand{\Ltotal}{\mathcal{L}}
\newcommand{\pos}{\mathrm{pos}}

\title{\textbf{Unitree Go2: System Dynamics, Output-Feedback Controller,\\
               and Reachable-Set Separation for Active Fault Diagnosis}}
\author{Technical Report}
\date{March 2026}

\begin{document}
\maketitle
\tableofcontents
\newpage

% =============================================================================
\section{Introduction}
% =============================================================================

This report describes the mathematical foundations and software implementation
of a reachable-set-based active fault diagnosis framework for the Unitree Go2
quadruped robot. The core idea is to design a \emph{separating input} or
\emph{output-feedback controller} that steers the robot along a trajectory
whose resulting reachable sets are \emph{disjoint} across different fault
hypotheses, enabling unambiguous fault identification from position
observations alone.

Three fault scenarios are considered:
\begin{enumerate}
  \item \textbf{Nominal}: the robot operates as designed.
  \item \textbf{Actuator Fault}: one or more actuators deliver reduced lateral
        and yaw authority ($\alpha, \beta \in [0.6, 0.8]$).
  \item \textbf{Sensor Fault}: the lateral-velocity sensor returns corrupted
        readings near zero, causing the dead-reckoned position estimate to
        drift from the true trajectory.
\end{enumerate}

For each scenario the reachable set of the robot's position is bounded using
\emph{interval arithmetic} and propagated forward in time via the natural
embedding (the tightest monotone inclusion function). A time-varying
output-feedback controller $u_k = u_k^{\rm nom} + K_k(y_k - y_k^{\rm nom})$
is then optimised so that the reachable sets of the three scenarios are
maximally separated at the final time horizon $T$. The separation condition
is the robot's analogue of \emph{collision avoidance} between probability
masses: the optimised controller prevents the position uncertainty tubes of
distinct fault modes from overlapping, guaranteeing that any observation of
the final position uniquely identifies the active fault.

% =============================================================================
\section{System Dynamics}
% =============================================================================

\subsection{Unicycle Model}

The Unitree Go2 is abstracted as a kinematic unicycle in the horizontal plane.
The state vector is
\[
  \xvec = \begin{bmatrix} p_x \\ p_y \\ \theta \end{bmatrix} \in \R^3,
\]
where $(p_x, p_y)$ is the robot's position in the world frame and $\theta$ is
its heading angle. The control input is
\[
  \uvec = \begin{bmatrix} v_x \\ v_y \\ \omega \end{bmatrix} \in \mathcal{U},
  \qquad
  \mathcal{U} = [-0.6,\,0.6]^3\;\text{m/s (or rad/s)},
\]
where $v_x$ and $v_y$ are body-frame forward and lateral velocities, and
$\omega$ is the commanded yaw rate.

The continuous-time equations of motion are
\begin{equation}
  \label{eq:unicycle}
  \begin{aligned}
    \dot{p}_x &= v_x \cos\theta - \alpha\, v_y \sin\theta, \\
    \dot{p}_y &= v_x \sin\theta + \alpha\, v_y \cos\theta, \\
    \dot{\theta} &= \beta\,\omega,
  \end{aligned}
\end{equation}
where the parameters $\alpha \in (0,1]$ and $\beta \in (0,1]$ encode actuator
effectiveness. Under nominal conditions $\alpha = \beta = 1$ and
\eqref{eq:unicycle} recovers the standard unicycle. An actuator fault reduces
both parameters uniformly.

\subsection{Nominal Trajectory}

For demonstration a circular reference trajectory of radius $r = 2.0\,\text{m}$
at angular rate $\omega_{\rm nom} = 0.3\,\text{rad/s}$ is used over a horizon
of $N = 20$ steps with $\Delta t = 0.5\,\text{s}$:
\[
  p_x^{\rm nom}(t) = r\cos(\omega_{\rm nom} t),
  \quad
  p_y^{\rm nom}(t) = r\sin(\omega_{\rm nom} t),
  \quad
  \theta^{\rm nom}(t) = \omega_{\rm nom} t + \tfrac{\pi}{2}.
\]
The corresponding nominal control is computed by differentiating the reference:
\[
  v_x^{\rm nom}(t) = -r\,\omega_{\rm nom}\sin(\omega_{\rm nom} t),
  \quad
  v_y^{\rm nom}(t) = r\,\omega_{\rm nom}\cos(\omega_{\rm nom} t),
  \quad
  \omega^{\rm nom}(t) = \omega_{\rm nom}.
\]

\begin{figure}[ht]
  \centering
  \includegraphics[width=\linewidth]{fig1_go2_nominal_trajectory.pdf}
  \caption{%
    \textbf{Left}: Nominal circular trajectory in the horizontal plane.
    Arrows indicate the robot heading $\theta$ at sampled time steps.
    \textbf{Right}: Nominal control inputs $v_x(t)$, $v_y(t)$, and $\omega(t)$.
  }
  \label{fig:nominal}
\end{figure}

% =============================================================================
\section{Fault Models}
% =============================================================================

\subsection{Actuator Fault}

An actuator fault reduces the robot's ability to track lateral and rotational
commands. Both the lateral gain $\alpha$ and the yaw gain $\beta$ drop below
their nominal values of 1:
\[
  \alpha \in [\alpha_{\rm lo},\,\alpha_{\rm hi}] = [0.60,\,0.80],
  \qquad
  \beta  \in [\beta_{\rm lo},\,\beta_{\rm hi}]   = [0.60,\,0.80].
\]
The dynamics remain \eqref{eq:unicycle} but with degraded parameters. As a
result, the trajectory drifts away from the nominal circle: lateral motion is
under-delivered and turns are shallower.

\subsection{Sensor Fault}

The lateral-velocity sensor returns a noisy reading near zero rather than the
true commanded $v_y$. The robot's onboard odometry integrates this corrupted
measurement, producing a \emph{dead-reckoned} position estimate $\hat{\xvec}$
that diverges from the true position. The corrupted velocity is modelled as
\begin{equation}
  \label{eq:sensor_fault}
  v_y^{\rm corrupted} = (1 - \omega^2)\,\alpha\,v_y
                        + \omega^2\,v_y^{\rm noise},
  \quad
  v_y^{\rm noise} \in [-\varepsilon,\,+\varepsilon],
  \quad \varepsilon = 0.25\;\text{m/s}.
\end{equation}
The weighting $(1-\omega^2)$ and $\omega^2$ smoothly interpolates between the
true measurement and the noise term as a function of the current yaw rate.
Under a pure sensor fault $\alpha = \beta = 1$; only the vy measurement is
corrupted. The dead-reckoned dynamics therefore are
\begin{equation}
  \label{eq:sensor_dyn}
  \begin{aligned}
    \dot{\hat{p}}_x &= v_x \cos\hat{\theta}
                       - v_y^{\rm corrupted}\sin\hat{\theta}, \\
    \dot{\hat{p}}_y &= v_x \sin\hat{\theta}
                       + v_y^{\rm corrupted}\cos\hat{\theta}, \\
    \dot{\hat{\theta}} &= \omega.
  \end{aligned}
\end{equation}

\begin{figure}[ht]
  \centering
  \includegraphics[width=\linewidth]{fig2_go2_fault_scenarios.pdf}
  \caption{%
    Visual summary of the three fault scenarios.
    \textbf{Left}: Nominal — full actuator authority and accurate sensors.
    \textbf{Centre}: Actuator fault — reduced lateral and yaw effectiveness
    ($\alpha, \beta \in [0.6, 0.8]$).
    \textbf{Right}: Sensor fault — corrupted $v_y$ measurement causes
    dead-reckoning drift.
  }
  \label{fig:faults}
\end{figure}

% =============================================================================
\section{Reachable-Set Propagation via Interval Arithmetic}
% =============================================================================

\subsection{Interval Arithmetic Primer}

An \emph{interval} $\Xivl = [\xvec^L, \xvec^U] \subset \R^n$ over-approximates
a set of states. For a function $f : \R^n \to \R^m$, an \emph{inclusion
function} $F$ satisfies $f(\Xivl) \subseteq F(\Xivl)$ component-wise. The
\emph{natural inclusion function} (natural embedding) replaces every real
operation with its interval counterpart and achieves the tightest possible
over-approximation when there is no variable repetition.

Tight bounds on $\cos\theta$ and $\sin\theta$ over $\theta \in [\theta^L,\theta^U]$
are computed by evaluating at the endpoints and checking whether any global
extrema (at multiples of $\pi/2$) lie within the interval:
\[
  \cos[\theta^L, \theta^U] \supseteq
  \bigl[\min(\cos\theta^L, \cos\theta^U,\,-\mathbf{1}_{[\pi]}),\,
        \max(\cos\theta^L, \cos\theta^U,\, \mathbf{1}_{[0]})\bigr],
\]
where $\mathbf{1}_{[\cdot]}$ equals 1 if $\theta \in [\theta^L, \theta^U]$
contains the indicated multiple of $\pi$, and 0 otherwise.

\subsection{Euler-Step Propagation}

For a fixed control input $\uvec_k$ and parameter interval $\Pivl$, the
state interval is advanced one Euler step:
\begin{equation}
  \label{eq:euler_ivl}
  \Xivl_{k+1} = \Xivl_k + \Delta t\cdot F(\Xivl_k, \uvec_k, \Pivl),
\end{equation}
where $F$ is the natural embedding of the right-hand side of
\eqref{eq:unicycle} (or \eqref{eq:sensor_dyn} for the sensor-fault scenario).
The interval arithmetic in $F$ handles the nonlinear $\cos\theta$ and
$\sin\theta$ terms and the bilinear multiplication $\alpha\,v_y$ exactly, so
$\Xivl_{k+1}$ is a guaranteed over-approximation of the true reachable set.

\subsection{Position Projection}

For fault diagnosis only the 2D position projection matters:
\[
  \pos(\Xivl) = [p_x^L, p_x^U] \times [p_y^L, p_y^U] \subset \R^2.
\]
The pairwise \emph{overlap volume} between two position intervals
$\pos(\Xivl^{(i)})$ and $\pos(\Xivl^{(j)})$ is
\begin{equation}
  \label{eq:overlap}
  \mathrm{Vol}\!\left(\pos(\Xivl^{(i)}) \cap \pos(\Xivl^{(j)})\right)
  = \prod_{d=1}^{2} \max\!\left(0,\;
    \min(p_d^{U,(i)}, p_d^{U,(j)}) - \max(p_d^{L,(i)}, p_d^{L,(j)})
  \right).
\end{equation}
When this quantity is zero the two position sets are \emph{separated}: no
observation of the robot's final position can be consistent with both fault
scenarios simultaneously, enabling unambiguous fault identification.

% =============================================================================
\section{Output-Feedback Controller}
% =============================================================================

\subsection{Controller Structure}

Rather than a constant open-loop excitation, a \emph{time-varying output-feedback}
controller is used that adapts to the current observation:
\begin{equation}
  \label{eq:controller}
  \uvec_k = \uvec_k^{\rm nom} + K_k\,(\yvec_k - \yvec_k^{\rm nom}),
\end{equation}
where
\begin{itemize}
  \item $\uvec_k^{\rm nom} \in \R^3$ is the nominal feedforward control computed
        from the reference trajectory,
  \item $K_k \in \R^{3\times m}$ is the time-varying feedback gain matrix
        optimised for fault separation,
  \item $\yvec_k = C\xvec_k$ is the observation with matrix $C \in \R^{m\times 3}$
        (identity for full-state observation),
  \item $\yvec_k^{\rm nom} = C\xvec_k^{\rm nom}$ is the nominal observation.
\end{itemize}

The controller \eqref{eq:controller} has two roles: the feedforward term tracks
the nominal trajectory, while the feedback term drives the robot to deviate
from the nominal in a way that is \emph{different} for each fault scenario,
thereby increasing the separation between their reachable sets.

\subsection{Interval Propagation under Feedback}

To propagate the reachable set forward under \eqref{eq:controller}, the
observation interval is first computed:
\[
  \Yvec_k = C\Xivl_k^{\rm meas} \oplus [-\delta, +\delta],
\]
where $\Xivl_k^{\rm meas}$ accounts for any sensor fault applied to the current
state interval and $\delta > 0$ is a uniform observation uncertainty. The
output-error interval is
\[
  \Delta\Yvec_k = \Yvec_k \ominus \yvec_k^{\rm nom}
  = [y^L_k - \yvec_k^{\rm nom},\; y^U_k - \yvec_k^{\rm nom}].
\]
The feedback control contribution is
\[
  \mathbf{U}_k^{\rm fb} =
  \left[
    \sum_j \left(K_{k,ij}\right)^+ \Delta y^L_{k,j}
    + \left(K_{k,ij}\right)^- \Delta y^U_{k,j},\;
    \sum_j \left(K_{k,ij}\right)^+ \Delta y^U_{k,j}
    + \left(K_{k,ij}\right)^- \Delta y^L_{k,j}
  \right]_{i=1}^{3},
\]
where $(\cdot)^+ = \max(\cdot,0)$ and $(\cdot)^- = \min(\cdot,0)$.
The total control interval is
$\mathbf{U}_k = \uvec_k^{\rm nom} + \mathbf{U}_k^{\rm fb}$.
The next state interval is then the convex hull of the two Euler steps
evaluated at the interval endpoints:
\begin{equation}
  \Xivl_{k+1} = \mathrm{conv}\!\left(
    \Xivl_k + \Delta t\,F(\Xivl_k,\uvec_k^{(L)},\Pivl),\;
    \Xivl_k + \Delta t\,F(\Xivl_k,\uvec_k^{(U)},\Pivl)
  \right),
\end{equation}
where $\uvec_k^{(L)}$ and $\uvec_k^{(U)}$ are the lower and upper bounds of
$\mathbf{U}_k$. This conservative union over-approximates the true reachable
set under all possible observations in $\Yvec_k$.

\subsection{Linear Output-Feedback (Immrax Formulation)}

An alternative, more elegant formulation encodes the controller directly into
the closed-loop system. The decision variable is $\theta = [\mathrm{vec}(K),\,
\mathbf{r}] \in \R^{12}$ (nine gain entries plus a three-dimensional
feedforward offset), and the closed-loop vector field is
\[
  f_{\rm cl}(\xvec, \theta, \pvec)
  = f\!\left(\xvec,\; \mathrm{clip}(K\xvec + \mathbf{r},\, \uvec^{\rm lo},
    \uvec^{\rm hi}),\; \pvec\right).
\]
Because $\theta$ is a concrete optimisation variable and only $\xvec$ and
$\pvec$ are abstract (intervals), the \texttt{immrax} natural embedding
correctly bounds $f_{\rm cl}$ over any box in state-parameter space for any
fixed $\theta$.

% =============================================================================
\section{Reachable-Set Separation (Avoiding Interval Overlap)}
% =============================================================================

\subsection{Separation as Collision Avoidance}

The core requirement for unambiguous fault diagnosis is that no observation
of the robot's final position $\yvec_T$ is simultaneously consistent with two
different fault scenarios. Formally, scenarios $i$ and $j$ are
\emph{distinguishable at time $T$} if and only if
\[
  \pos\!\left(\Xivl_T^{(i)}\right) \cap \pos\!\left(\Xivl_T^{(j)}\right)
  = \emptyset.
\]
This is precisely a \emph{collision avoidance} condition between two sets
moving in 2D position space: the interval reachable sets must not intersect.
The total volume of all pairwise intersections serves as the optimisation
objective, which the feedback controller drives to zero.

\subsection{Separation Loss}

Let $\mathcal{S} = \{1, 2, 3\}$ index the three scenarios. Define
\begin{equation}
  \label{eq:loss_sep}
  \Lsep(K) = \sum_{\substack{i,j \in \mathcal{S} \\ i < j}}
    \mathrm{Vol}\!\left(
      \pos\!\left(\Xivl_T^{(i)}(K)\right) \cap
      \pos\!\left(\Xivl_T^{(j)}(K)\right)
    \right),
\end{equation}
where $\Xivl_T^{(i)}(K)$ is the reachable set at time $T$ under controller
$K$ and fault parameters $\pvec^{(i)}$.

\subsection{Regularised Objective}

To prevent degenerate solutions with arbitrarily large gains, a Frobenius-norm
regulariser is added:
\begin{equation}
  \label{eq:loss_total}
  \Ltotal(K) = \lambda_{\rm sep}\,\Lsep(K)
             + \lambda_{\rm track}\,\|K\|_F^2,
  \qquad
  \lambda_{\rm sep} = 1.0,\quad
  \lambda_{\rm track} = 0.01.
\end{equation}
The regularisation term keeps the feedback gains small, ensuring the actual
trajectory remains close to the nominal circular reference.

% =============================================================================
\section{Optimisation}
% =============================================================================

\subsection{Gradient Descent}

Because $\Ltotal$ is differentiable with respect to $K$ (the interval Euler
step is a composition of smooth JAX operations), the gains are optimised by
gradient descent:
\begin{align*}
  K^{(\ell+1)} &= K^{(\ell)} - \eta\,\nabla_K\,\Ltotal(K^{(\ell)}),
  \quad \ell = 0, 1, \ldots, L_{\rm iter}-1,
\end{align*}
with learning rate $\eta = 10^{-3}$ and $L_{\rm iter} = 100$ iterations.
The gradient $\nabla_K\,\Ltotal$ is computed by automatic differentiation
(\texttt{jax.grad}) through the interval propagation.

\subsection{Multi-Start GPU Parallelism (Linear Feedback)}

For the linear output-feedback formulation ($\theta \in \R^{12}$) a
\emph{multi-start} strategy is employed: $R = 100$ random initialisations of
$\theta$ are optimised \emph{in parallel} using \texttt{jax.vmap} over a
\texttt{jax.lax.fori\_loop}, and the restart with the lowest final loss is
returned. Each restart is projected at every step onto the feasible box
$K_{ij} \in [-K_{\rm max}, K_{\rm max}]$, $\mathbf{r} \in \mathcal{U}$.

\begin{center}
\fbox{\begin{minipage}{0.82\textwidth}
\textbf{Algorithm: Multi-Start Output-Feedback Optimisation}\\[4pt]
\textbf{Input}: $\Xivl_0$, $\{\Pivl^{(i)}\}$, step $\Delta t$, horizon $N$,
restarts $R$, learning rate $\eta$, iterations $L_{\rm iter}$\\[2pt]
1.\ Sample $\{\theta_0^{(r)}\}_{r=1}^R \sim \mathcal{N}(\bar\theta,\,0.01\,I_{12})$\\
2.\ \textbf{for} $\ell = 1,\ldots, L_{\rm iter}$
    \quad\textit{(all $r$ in parallel via \texttt{jax.vmap})}:\\
\quad\quad $g^{(r)} \leftarrow \nabla_\theta\,\Lsep(\theta^{(r)})$
           \quad for all $r$\\
\quad\quad $\theta^{(r)} \leftarrow \mathrm{clip}
           \!\left(\theta^{(r)} - \eta\,g^{(r)},\;
           \Theta_{\rm lo},\,\Theta_{\rm hi}\right)$\\
3.\ $r^* \leftarrow \arg\min_r \Lsep(\theta^{(r)})$\\
\textbf{Return}: $\theta^{(r^*)}$
\end{minipage}}
\end{center}

% =============================================================================
\section{Simulation Results}
% =============================================================================

\subsection{Reachable Sets}

Figure~\ref{fig:openloop} shows the position reachable sets for all three
fault scenarios under the nominal open-loop control $u = u^{\rm nom}$.
At $t = T = 10\,\text{s}$ the three sets overlap with total pairwise volume
$0.298\,\text{m}^2$: a robot observed at any of those positions could belong
to any of the three scenarios. Fault diagnosis is therefore impossible.

\begin{figure}[ht]
  \centering
  \includegraphics[width=0.7\linewidth]{fig3_go2_reachable_sets_openloop.pdf}
  \caption{%
    Position reachable sets for Nominal (blue), Actuator Fault (green),
    and Sensor Fault (red) under the open-loop controller ($K=0$).
    Rectangles show the interval hull at each time step; the dashed line
    is the nominal circular reference. The sets overlap heavily at $t=T$,
    making the scenarios indistinguishable.
  }
  \label{fig:openloop}
\end{figure}

Figure~\ref{fig:feedback} shows the same reachable sets after optimising the
output-feedback gains $K$. The separation loss drops from $0.298\,\text{m}^2$
to $0.114\,\text{m}^2$ — a $61.7\%$ reduction — confirming that the optimised
feedback steers the three scenario tubes onto substantially different regions
of position space. With further iterations or a higher-dimensional excitation
the overlap can be driven to zero.

\begin{figure}[ht]
  \centering
  \includegraphics[width=0.7\linewidth]{fig4_go2_reachable_sets_feedback.pdf}
  \caption{%
    Position reachable sets under the optimised output-feedback controller.
    The three scenario tubes are steered apart by the feedback term $K_k
    (y_k - y_k^{\rm nom})$, minimising overlap at $t=T$.
  }
  \label{fig:feedback}
\end{figure}

\subsection{Feedback Gains}

Figure~\ref{fig:gains} shows all nine entries of the time-varying gain matrix
$K[t] \in \R^{3\times 3}$. The gains are non-trivial primarily in the middle
of the trajectory where the nonlinear coupling between heading $\theta$ and
lateral velocity creates the most leverage for separating the fault-induced
trajectories. Gains are small in magnitude (bounded by $\lambda_{\rm track}$
regularisation), confirming that the feedback correction is a perturbation
around the nominal feedforward trajectory rather than a fundamentally
different control policy.

\begin{figure}[ht]
  \centering
  \includegraphics[width=\linewidth]{fig5_go2_feedback_gains.pdf}
  \caption{%
    Time-varying feedback gain matrix $K[t] \in \R^{3\times 3}$ over the
    20-step horizon. Each sub-plot shows how aggressively one control
    channel ($v_x$, $v_y$, $\omega$) responds to a unit error in one
    observation channel ($p_x$, $p_y$, $\theta$).
  }
  \label{fig:gains}
\end{figure}

\subsection{Overlap Comparison}

Figure~\ref{fig:overlap} quantifies the reduction in pairwise reachable-set
overlap achieved by the optimised feedback controller relative to the open-loop
baseline. The overlap between the \emph{Nominal} and \emph{Actuator Fault}
scenarios is most significantly reduced because the feedback gain directly
amplifies the rotational divergence induced by the reduced yaw gain $\beta$.
The \emph{Sensor Fault} scenario is inherently harder to separate because the
fault only affects the dead-reckoned estimate (not the physical trajectory),
so positional divergence accumulates more slowly.

\begin{figure}[ht]
  \centering
  \includegraphics[width=0.85\linewidth]{fig6_go2_overlap_comparison.pdf}
  \caption{%
    Pairwise position-interval overlap at $t=T$ for the open-loop ($K=0$)
    and optimised feedback controllers. Lower values indicate better fault
    distinguishability; zero means the scenarios are fully separated.
  }
  \label{fig:overlap}
\end{figure}

% =============================================================================
\section{Key Implementation Details}
% =============================================================================

\paragraph{Interval tightness.}
The natural embedding achieves the tightest possible interval bound when each
variable appears at most once in the expression. For the unicycle dynamics
\eqref{eq:unicycle}, the state $\theta$ appears in both $\cos\theta$ and
$\sin\theta$, so dependency effects are minimal and the bounds are tight in
practice.

\paragraph{Sensor-fault model.}
The blending factor $\omega^2$ in \eqref{eq:sensor_fault} ensures that the
sensor noise has no effect when the robot moves straight ($\omega = 0$) and
full effect during rapid turns. This captures the physical intuition that a
yaw-rate-dependent sensor corruption accumulates angular error primarily during
rotation.

\paragraph{Control projection.}
All control inputs are projected onto the feasible box $\mathcal{U} =
[-0.6,0.6]^3$ after every gradient step. For the linear feedback formulation,
the gain matrix entries are also clipped to $[-K_{\rm max}, K_{\rm max}]$ with
$K_{\rm max} = 2.0$.

\paragraph{JAX implementation.}
All interval propagation and gradient computation are implemented in JAX,
enabling JIT compilation and GPU acceleration. The multi-start parallelism
is achieved with a single \texttt{jax.vmap} wrapping the per-restart loss
function, giving wall-clock speedup proportional to the number of available
GPU cores.

% =============================================================================
\section{Summary}
% =============================================================================

Table~\ref{tab:summary} summarises the key parameters of the system and
optimisation problem.

\begin{table}[ht]
\centering
\caption{Summary of system and optimisation parameters.}
\label{tab:summary}
\begin{tabular}{@{}lll@{}}
\toprule
\textbf{Quantity} & \textbf{Symbol} & \textbf{Value} \\
\midrule
State dimension         & $n$               & 3 ($p_x, p_y, \theta$) \\
Control dimension       & $m$               & 3 ($v_x, v_y, \omega$) \\
Observation dimension   & $p$               & 3 (full state) \\
Control bounds          & $\mathcal{U}$     & $[-0.6,\,0.6]^3$ \\
Time-step               & $\Delta t$        & 0.5 s \\
Horizon                 & $T = N\Delta t$   & 10 s ($N=20$) \\
Actuator fault range    & $\alpha,\beta$    & $[0.60,\,0.80]$ \\
Sensor noise bound      & $\varepsilon$     & 0.25 m/s \\
Observation uncertainty & $\delta$          & 0.05 \\
Separation weight       & $\lambda_{\rm sep}$   & 1.0 \\
Regularisation weight   & $\lambda_{\rm track}$ & 0.01 \\
Learning rate           & $\eta$            & $10^{-3}$ \\
Iterations              & $L_{\rm iter}$    & 100 \\
Gain matrix bound       & $K_{\rm max}$     & 2.0 \\
\bottomrule
\end{tabular}
\end{table}

The optimised output-feedback controller achieves reliable separation of the
three fault scenarios by exploiting the state-dependent nonlinearity of the
unicycle model. The feedforward term tracks the nominal circular trajectory
while the time-varying feedback term steers each scenario's reachable set onto
a distinct region of position space, enabling fault identification from a
single position observation at time $T$.

\end{document}
